\section{Rank and the basic tensor}\label{sec:foundations}

Fix an arbitrary field $\F$.  All tensor identities and linear maps below
are over $\F$, and all logarithms are natural.  A tensor is viewed as a
trilinear form in three disjoint
sets of variables.  Its rank $\R(T)$ is the least number of products of
three linear forms, one on each side, whose sum is $T$.  A
\emph{restriction} substitutes linear forms separately on the three
sides.  Restrictions cannot increase rank.  A direct sum uses disjoint
variables, and a tensor product pairs the variable indices on each side;
rank is subadditive under direct sum and submultiplicative under tensor
product.  These statements follow by applying the indicated operations
to rank decompositions.

For positive integers $a,b,c$, write
\[
 \langle a,b,c\rangle
   =\sum_{i=1}^{a}\sum_{j=1}^{b}\sum_{k=1}^{c}
       x_{ij}y_{jk}z_{ki}.
\]
This is the tensor for multiplying an $a\times b$ matrix by a
$b\times c$ matrix.  Its \emph{volume} is $abc$.  Pairing the three
matrix indices gives
\[
 \langle a,b,c\rangle\otimes\langle a',b',c'\rangle
     \cong\langle aa',bb',cc'\rangle.
\]
Cyclically permuting the three variable spaces replaces the parameters
by $(b,c,a)$ or $(c,a,b)$ and preserves rank.

\subsection{Exact rank and the arithmetic exponent}

Put $R_d=\R(\langle d,d,d\rangle)$ and
\begin{equation}\label{eq:rank-exponent}
 w=\inf_{d\geq2}\frac{\log R_d}{\log d},
 \qquad d\text{ an integer}.
\end{equation}
The flattening that sends a $z$ coordinate to its coefficient bilinear
form has rank $d^2$: the $d^2$ forms
$\sum_jx_{ij}y_{jk}$ are linearly independent.  A rank-one tensor has
flattening rank at most one.  Thus
$d^2\leq R_d\leq d^3$ and $2\leq w\leq3$.  In particular,
\begin{equation}\label{eq:square-rank-lower}
 R_d\geq d^w\qquad(d\geq1),
\end{equation}
where $R_1=1$ handles the unit case.

Let $\omega$ denote the arithmetic exponent of square matrix
multiplication over $\F$, counting scalar additions and multiplications.
A fixed rank-$R_d$ decomposition gives a bilinear matrix multiplication
identity with $R_d$ products.  This identity remains valid on matrix
blocks: every product has its $x$-linear factor before its $y$-linear
factor, and comparison of coefficients proves the same identity for
noncommuting blocks with central coefficients in $\F$.  At size $d^k$
the resulting arithmetic cost satisfies
\[
 C_k\leq R_d C_{k-1}+O_d(d^{2k}).
\]
The implied constant may depend on the chosen decomposition.  Since
$R_d\geq d^2$, this gives $C_k=O_d(kR_d^k)$, including the boundary
case $R_d=d^2$.  Padding an arbitrary size to the next power of $d$
therefore gives cost $O_d(n^{\log_d R_d}\log(n+1))$.  The logarithmic
factor is absorbed by $n^\eta$ for every $\eta>0$; taking the
infimum in \eqref{eq:rank-exponent} proves
\begin{equation}\label{eq:omega-rank-exponent}
 \omega\leq w.
\end{equation}

We also need a uniform upper bound on exact square rank.  Given
$\varepsilon>0$, choose one integer $d\geq2$ with
$\theta=\log_d R_d<w+\varepsilon$.  For $h\geq1$, tensor powers and
padding give
\begin{equation}\label{eq:square-rank-upper}
 \R(\langle h,h,h\rangle)
 \leq R_d^{\lceil\log_d h\rceil}
 \leq C_\varepsilon h^{w+\varepsilon},
 \qquad C_\varepsilon=R_d.
\end{equation}
This argument uses one approximating dimension; it does not require the
infimum defining $w$ to be attained.

We use the following identical-summand form of Sch\"onhage's asymptotic
sum inequality~\cite{Schonhage1981}, with a proof from exact rank.

\begin{lemma}[Direct sums of identical matrix tensors]\label{lem:direct-sum}
Let $s,a,b,c$ be positive integers.  If
\[
 \R\left(\bigoplus_{j=1}^{s}\langle a,b,c\rangle\right)\leq R_*,
\]
then
\begin{equation}\label{eq:identical-direct-sum}
 s(abc)^{w/3}\leq R_*.
\end{equation}
\end{lemma}

\begin{proof}
First consider $s$ copies of $\langle d,d,d\rangle$.  When $s=1$,
the assertion $sd^w\leq R_*$ is \eqref{eq:square-rank-lower}.
When $d=1$, the direct sum is $\sum_{j=1}^s x_jy_jz_j$, whose
flattening has rank $s$, so the assertion also holds.  We may therefore
assume $s,d\geq2$.

Fix $\varepsilon>0$ and the constant in
\eqref{eq:square-rank-upper}.  For all sufficiently large integers $k$,
set
\[
 h_k=\left\lfloor
       (s^k/C_\varepsilon)^{1/(w+\varepsilon)}
       \right\rfloor.
\]
Then $h_k\geq2$,
$\R(\langle h_k,h_k,h_k\rangle)\leq s^k$, and
\[
 \lim_{k\to\infty}\frac{\log h_k}{k}
      =\frac{\log s}{w+\varepsilon}.
\]
The $k$th tensor power of the given direct sum is a direct sum of
$s^k$ copies of $\langle d^k,d^k,d^k\rangle$, of rank at most $R_*^k$.
A rank-$r$ decomposition of $\langle h_k,h_k,h_k\rangle$, with
$r\leq s^k$, expresses it as a restriction of $r$ independent scalar
products.  Tensor this restriction with
$\langle d^k,d^k,d^k\rangle$ and discard unused copies.  It gives
$\langle d^kh_k,d^kh_k,d^kh_k\rangle$ as a restriction of the
available direct sum.  Consequently
\[
 R_*^k\geq(d^kh_k)^w.
\]
Taking logarithms, dividing by $k$, and then taking $k\to\infty$
at fixed $\varepsilon$ yields
\[
 \log R_*\geq w\log d+\frac{w}{w+\varepsilon}\log s.
\]
Letting $\varepsilon\downarrow0$ proves $sd^w\leq R_*$.

For the rectangular case, let
$U=\bigoplus_{j=1}^s\langle a,b,c\rangle$ and tensor $U$ with its
two cyclic rotations.  Distributivity gives $s^3$ independent copies of
\[
 \langle a,b,c\rangle\otimes\langle b,c,a\rangle
       \otimes\langle c,a,b\rangle
       \cong\langle abc,abc,abc\rangle,
\]
and the product has rank at most $R_*^3$.  The square case gives
$s^3(abc)^w\leq R_*^3$, whose cube root is
\eqref{eq:identical-direct-sum}.  The same proof covers $abc=1$
through the unit square case.
\end{proof}

The hypothesis in Lemma~\ref{lem:direct-sum} specifies one common
\emph{oriented} triple $(a,b,c)$ for all summands.  In the construction
below, exact history counts and common terminal type restrictions will
provide these identical dimensions.

\subsection{An integer coefficient construction}

The relation between approximate bilinear identities and exact algorithms
was studied by Bini~\cite{Bini1980}.  Here a direct coefficient calculation
gives an exact rank bound over the original field, including when that
field is finite.

For an integer $q\geq1$, define the symmetric Coppersmith--Winograd
tensor~\cite{CoppersmithWinograd1990}
\begin{align*}
 \CW_q={}&\sum_{i=1}^q
  (x_0y_iz_i+x_iy_0z_i+x_iy_iz_0)\\
 &+x_{q+1}y_0z_0+x_0y_{q+1}z_0+x_0y_0z_{q+1}.
\end{align*}
We use $q=5$.  The following exact-rank estimate makes the field
independence and the polynomial overhead explicit.

\begin{lemma}[Rank bound for powers of the basic tensor]\label{lem:cw-rank}
Over every field and for every integer $m\geq0$,
\begin{equation}\label{eq:cw-power-rank}
 \R(\CW_q^{\otimes m})
      \leq(q+2)^m\binom{3m+2}{2}.
\end{equation}
\end{lemma}

\begin{proof}
We use the Laurent-polynomial identity of Coppersmith and
Winograd~\cite[Section~7, equation~(10)]{CoppersmithWinograd1990}
over the integers, with a formal variable $t$. Set
$X_1=\sum_{i=1}^q x_i$, $Y_1=\sum_{i=1}^q y_i$, and
$Z_1=\sum_{i=1}^q z_i$, and consider
\begin{align}
 P(t)={}&t^{-2}\sum_{i=1}^q
        (x_0+tx_i)(y_0+ty_i)(z_0+tz_i)\notag\\
 &-t^{-3}(x_0+t^2X_1)(y_0+t^2Y_1)(z_0+t^2Z_1)\notag\\
 &+(t^{-3}-qt^{-2})
        (x_0+t^3x_{q+1})(y_0+t^3y_{q+1})(z_0+t^3z_{q+1}).
        \label{eq:cw-laurent}
\end{align}
The only possible negative degrees are $-3,-2,-1$.  Their
coefficients are, respectively,
\[
 (-1+1)x_0y_0z_0,\qquad
 (q-q)x_0y_0z_0,\qquad
 (1-1)(X_1y_0z_0+x_0Y_1z_0+x_0y_0Z_1).
\]
The degree-zero coefficient is precisely $\CW_q$.  Thus
$P(t)=\CW_q+tQ(t)$ for a polynomial $Q$ with integer coefficients.
In particular,
\begin{equation}\label{eq:cw-constant-power}
 [t^0]P(t)^{\otimes m}=\CW_q^{\otimes m}.
\end{equation}
These coefficient identities remain valid after reduction to any field.

Regard \eqref{eq:cw-laurent} as $q+2$ products of three
Laurent-polynomial linear forms, absorbing the scalar prefactor of
each product into its $x$ form.  In one tensor position the possible
degrees on the $x$ side lie, respectively, in
\[
 \{-2,-1\},\qquad\{-3,-1\},\qquad\{-3,-2,0,1\}.
\]
The last set includes both terms of $t^{-3}-qt^{-2}$ within a single
linear form.  On each of the other two sides the degrees lie between
$0$ and $3$.

Expanding $P(t)^{\otimes m}$ therefore gives $(q+2)^m$ products,
each of the form $A(t)B(t)C(t)$ with
\[
 A(t)=\sum_{a=-3m}^{m}A_at^a,\qquad
 B(t)=\sum_{b=0}^{3m}B_bt^b,\qquad
 C(t)=\sum_{c=0}^{3m}C_ct^c.
\]
Here each coefficient is one linear form in the corresponding
tensor-product variables.  The constant coefficient of this product
is
\[
 \sum_{\substack{b,c\geq0\\b+c\leq3m}}
       A_{-b-c}B_bC_c,
\]
a sum of at most $\binom{3m+2}{2}$ rank-one tensors.  Summing over
the $(q+2)^m$ choices and using \eqref{eq:cw-constant-power} proves
\eqref{eq:cw-power-rank}.  For $m=0$, both sides equal $1$ under
the empty tensor product convention.  Only coefficient addition and
multiplication were used; no interpolation or field division is needed.
\end{proof}

\subsection{Strands and shapes}

Give the label $0$ weight $0$, the labels $1,\ldots,q$ weight $1$,
and the label $q+1$ weight $2$.  Every term of $\CW_q$ has total
weight $2$.  A \emph{strand of length $\ell$} is a copy of
$\CW_q^{\otimes\ell}$; its variable indices are words of length
$\ell$.  For a triple $g=(g_0,g_1,g_2)$ of nonnegative integers
with $g_0+g_1+g_2=2\ell$, let $T_{\ell,g}$ be the restriction in
which the weight sums on the three sides are $g_0,g_1,g_2$.
The coordinates of $g$ refer to sides, whereas the weights $0,1,2$
refer to individual labels.  A shape is \emph{positive} if all three
coordinates are positive.  When $\ell$ is even, splitting a strand
into its two consecutive halves gives child shapes $u$ and $g-u$,
where $0\leq u\leq g$ coordinatewise and $\sum_i u_i=\ell$.

A zero coordinate gives an especially simple terminal tensor.
Suppose $g_0=0$.  The $x$ word is then forced to be the all-zero
word.  In each position a label on the $y$ side uniquely determines
the $z$ label: $0$ and $q+1$ are exchanged, and every label
$i\in\{1,\ldots,q\}$ is fixed.  Thus the other two sides are
matched bijectively, and $T_{\ell,g}$ is
$\langle1,1,M\rangle$, where
\[
 M=[z^{g_1}](1+qz+z^2)^\ell.
\]
For $g_1=0$ or $g_2=0$ the corresponding orientations are
$\langle M,1,1\rangle$ or $\langle1,M,1\rangle$.  Restricting
the matched words by corresponding type conditions on their two
sides preserves this form, with $M$ replaced by the number of
allowed matched pairs; an empty matching gives the zero tensor.
This observation will supply the terminal volume counts.
